\subsection{Theory}
Watts-Strogatz graphs are a type of network that are constructed as follows: 
\begin{itemize}
    \item $N$ nodes are arranged on a ring
    \item In the beginning all nodes are connected to their $k/2$ closest neighbours to both sides (every node has degree $k$)
    \item Each edge is rewired with a probability of $p$ (no self-loops or duplicate edges allowed)
\end{itemize}

Independent of $p$, all graphs with the same $N$ and $k$ will have the same number of edges.
However, these edges will be placed different depending on $p$ as a rewired edge can be placed anywhere and not only between nodes situated close to each other. 
For small $p\ne0$ the rewiring will lead to a few long-range shortcuts. For $p=1$ all edges are rewired and we get a completely random graph. Below is a visualisation of the graph.

\begin{figure}[H]
    \centering
    \includegraphics[width=0.5\linewidth]{figs/WSgraph_visual.pdf}
    \caption{A Watts-Strogatz graph with $N=30$, $k=6$ and $p=0.1$. The lines inside the circle are the random edges generated with a non-zero $p$.}
    \label{fig:WScircleplot}
\end{figure}

\subsubsection{Degrees and degree distribution}


For a graph node $i$ the degree $d_i$ describes the number of edges incident to that node. It is also possible to calculate the average degree for the whole graph using 

\begin{equation}
    \bar{d} = \frac{1}{N} \sum^N_{i = 1} d_i 
\end{equation}

where N is the total number of nodes. In the case of a Watts-Strogatz graph the average degree is equal to $k$ as initially every node is connected to $k$ other nodes. At $p=0$ there are no random connections added and every node has exactly $k$ connections and so for $p=0$ every node has the degree is $d_i = k$. As $p$ is increased, random connections appear and thus $d_i$ starts to vary between nodes but the mean $d_i$ would still be equal to $k$.

\subsubsection{Connected components}

An undirected graph is called connected if every node can be connected to every other node by a path. If a graph is disconnected then a connected subgraph of the original graph is called a connected component. A connected component can be studied by computing different distance quantities for it such as the shortest path length, the radius and the diameter.



\subsubsection{Shortest distance, diameter and average path length}

In a Watts-Strogatz small world graph, edges connect two arbitrary nodes $i$ and $j$ through a path, a sequence of edges. A graph is connected if any $i$ can be reached from any $j$ through any path of any length. With the assumption that every graph is connected, define the diameter and average path length. Begin with dist($i, j$), which is the shortest-path distance between $i$ and $j$, i.e. the least number of edges which connect the two nodes. The largest possible shortest path between two nodes in the graph is the diameter, written as 
\begin{align}
    \text{diam(G)} = \max_{i,j} \text{dist}(i,j).
    \label{eq:WSdiameter}
\end{align}

The average path length is instead the average shortest path length between $i$ and $j$ and is defined as follows:
\begin{align}
    \ell(\text{G}) = \frac{1}{N(N-1)}\sum_{i\neq j}\text{dist}(i,j)
    \label{eq:WSavgshortestpath}
\end{align}

A graph can also be disconnected if there are nodes or sets of nodes that cannot be reached from any other node. The diameter and average path length would still be defined in the same way as for a connected graph, but the numerical value would be different once computed.


\subsubsection{Clustering Coefficient}
The local clustering coefficient measures how many neighbouring nodes of one node are connected among themselves. 
Using the example of friendship, it describes how many of your friends are friends with your other friends.

Mathematically, the local clustering coefficient is the ratio of the number of edges between the neighbours and the possible number of edges between the neighbours.
Thus, a clustering coefficient of 1 means high clustering, all neighbours are connected, while 0 means no clustering or none of a node's neighbours are connected.
Let $d_i$ denote the degree, or number of neighbours, of node $i$. Then the maximum number of connections between node $i$'s neighbours is
\begin{equation*}
	\frac{d_i(d_i-1)}{2}
\end{equation*}

With $e_i$ denoting the number of these edges that are present, the local clustering coefficient becomes 
\begin{equation*}
C_i = \frac{2e_i}{d_i(d_i-1)}
\end{equation*}
If a node has only one or no neighbour ($d_i < 2$), the local clustering coefficient is automatically set to $C_i=0$.

A more general number than the local clustering coefficient is the average clustering coefficient, which averages over all nodes of a graph. 
\begin{equation*}
C(G) = \frac{1}{N} \sum_{i=1}^{N}C_i
\end{equation*}
where $G$ denotes the graph and $N$ is the number of nodes. This number will be calculated for all graphs and plotted together with the average path length. 




